\documentclass[11pt]{article}
\usepackage{wok-editorial}

\begin{document}

\woktitle{Concurso de Contos (Simulação e Algoritmos Gulosos)}{Mundo do Código}{\today}{Algoritmo Guloso, Simulação, Lógica Básica, Strings}

\begin{objetivobox}
\begin{itemize}
    \item \textbf{Paradigma Guloso (Greedy):} Compreender e aplicar a escolha ótima local (alocar o máximo possível de palavras na linha atual) para atingir o ótimo global (minimizar o número de páginas).
    \item \textbf{Modelagem de Estado Discreto:} Desenvolver a capacidade de rastrear e atualizar múltiplas variáveis dependentes (caracteres, linhas, páginas) através de uma simulação passo-a-passo.
    \item \textbf{Tratamento de Casos de Borda:} Gerenciar adequadamente regras condicionais específicas, como o não-cômputo de espaços no início de linhas e transições exatas nos limites de capacidade.
    \item \textbf{Eficiência de I/O em Strings:} Analisar estratégias de leitura sob demanda (streaming) em contraposição ao carregamento completo de dados na memória.
\end{itemize}
\end{objetivobox}

\vspace{0.5em}

\section{Disciplinas Relacionadas}
\begin{table}[h]
\centering
\begin{tabularx}{\textwidth}{l X}
\toprule
\textbf{Disciplina} & \textbf{Tópicos Relevantes} \\
\midrule
Algoritmos e Estruturas de Dados & Paradigmas de projetos de algoritmos (Algoritmos Gulosos), Análise Assintótica. \\
Lógica de Programação & Estruturas de controle de fluxo, acumulação de estado, laços de repetição, simulação de processos. \\
Sistemas Operacionais & Gerenciamento de buffers, leitura de fluxos de entrada (streaming) vs. alocação contígua em memória. \\
\bottomrule
\end{tabularx}
\caption{Mapeamento para currículos acadêmicos de Ciência da Computação.}
\end{table}

\vspace{1em}

\tagcloud{Algoritmo Guloso, Simulação, Strings, Variáveis de Estado, Processamento de Texto, Formatação, Laços de Repetição, Otimização de I/O}

\newpage

\section{Editorial Completo}

O problema ``Concurso de Contos'' é um clássico de simulação baseada em restrições de formatação. O objetivo é determinar o número mínimo de páginas necessárias para imprimir um texto contínuo, dado um número predeterminado de caracteres por linha e linhas por página.

\subsection{A Natureza do Problema: Formatação de Texto}

Diferente de problemas avançados de diagramação (como o \textit{Text Justification} ou o algoritmo do \TeX, que utilizam Programação Dinâmica para minimizar a ``feiura'' das quebras de linha), este problema não se importa com a estética do alinhamento direito. A única restrição aqui é o limite rígido superior de caracteres por linha, garantindo-se que cada palavra e seu respectivo espaço couberem no espaço restante da linha atual.

Isto simplifica a natureza do problema e permite uma abordagem direta, conhecida como \textbf{Algoritmo Guloso} (\textit{Greedy Algorithm}).

\subsection{O Paradigma Guloso}

Um algoritmo guloso constrói uma solução global fazendo escolhas localmente ótimas a cada etapa. No contexto da diagramação de texto, a escolha gulosa se expressa da seguinte forma: 
\textit{``Se uma palavra couber na linha atual, ela deve ser colocada nela. Caso contrário, e somente neste caso, move-se para a próxima linha.''}

\begin{notebox}[Por que a escolha gulosa é ótima neste caso?]
Poderia haver algum cenário em que mover uma palavra para a linha de baixo precocemente (mesmo havendo espaço na linha atual) resultaria em um número menor de páginas no final?
A resposta matemática é não. Deixar espaço livre em uma linha sem necessidade apenas desloca todas as palavras seguintes para a frente, o que pode forçar novas quebras de linha mais adiante, potencialmente aumentando o número de páginas. Não existe penalidade para linhas com larguras variadas. Assim, maximizar a ocupação de caracteres em cada linha sempre minimizará o total de linhas usadas e, consequentemente, o número total de páginas.
\end{notebox}

\subsection{Modelagem Matemática e Máquina de Estados}

Para simular o processo, modelaremos o problema como uma máquina de estados finitos simples que atualiza suas variáveis a cada entrada de palavra.

Definimos o estado pelo triplete $(P_{atual}, L_{atual}, C_{atual})$, onde:
\begin{itemize}
    \item $P_{atual}$: Número de páginas usadas até o momento (inicialmente 1).
    \item $L_{atual}$: Número de linhas ocupadas na página atual (inicialmente 1).
    \item $C_{atual}$: Quantidade de caracteres já escritos na linha atual (inicialmente 0).
\end{itemize}

Além disso, sejam $C_{max}$ a quantidade máxima de caracteres por linha, e $L_{max}$ a quantidade máxima de linhas por página. Considere uma sequência de palavras de comprimentos $w_1, w_2, \ldots, w_N$.

Para o processamento da $i$-ésima palavra de comprimento $w_i$, calculamos a quantidade de espaço $S_i$ que esta palavra exigirá na linha atual.
Se for a primeira palavra a ser impressa na linha corrente ($C_{atual} = 0$), nenhum caractere de espaço é necessário antes dela. Se já houverem palavras na linha corrente ($C_{atual} > 0$), é obrigatório adicionar 1 caractere em branco como separador. 

A regra pode ser expressa como:
$$
S_i = \begin{cases} 
w_i, & \text{se } C_{atual} = 0 \\
w_i + 1, & \text{se } C_{atual} > 0 
\end{cases}
$$

A transição de estados se dá verificando se $C_{atual} + S_i \leq C_{max}$.

\textbf{Caso 1: A palavra cabe na linha atual.}
O estado é atualizado adicionando a exigência ao acumulador de caracteres da linha atual:
$$C_{atual} \leftarrow C_{atual} + S_i$$

\textbf{Caso 2: A palavra \underline{não} cabe na linha atual.}
Um transbordamento ocorre. A linha deve ser finalizada.
$$L_{atual} \leftarrow L_{atual} + 1$$
$$C_{atual} \leftarrow w_i$$
Note que na nova linha, a palavra se torna a primeira, portanto, ocupará apenas $w_i$ (sem espaço precedente).

\begin{warnbox}[Cascata de Eventos: Checagem de Página]
Imediatamente após incrementar a linha atual ($L_{atual}$), surge uma ramificação lógica crítica: devemos verificar se o novo valor de $L_{atual}$ excede o limite estrito da página.
Se $L_{atual} > L_{max}$, então o estado transita para:
$$P_{atual} \leftarrow P_{atual} + 1$$
$$L_{atual} \leftarrow 1$$
A omissão dessa cascata é a principal fonte de falhas algorítmicas (Wrong Answer) neste desafio.
\end{warnbox}

\subsection{Estratégias de Leitura e Otimização de I/O}

No enunciado, as palavras são descritas como texto integral, separadas por espaços em branco. Do ponto de vista da arquitetura de computadores, existem duas formas primárias de lidar com essa entrada:

1. \textbf{Materialização na Memória (\textit{Buffered Read}):} Consiste em ler toda a sequência de caracteres de uma vez (por exemplo, usando \texttt{readline} até o fim), armazenar em uma string colossal na memória, aplicar uma função de particionamento (como \texttt{split}) gerando um grande arranjo de strings e, finalmente, iterar sobre o array calculando os comprimentos. Essa abordagem tem complexidade de espaço $\mathcal{O}(S_{total})$ e pode acionar instabilidades de coleta de lixo (\textit{Garbage Collection}) se a string for massiva.

2. \textbf{Processamento de Fluxo Contínuo (\textit{Stream Processing}):} Lê as palavras sob demanda, ignorando ativamente os separadores, tratando e descartando cada palavra em sequências mínimas da memória. O estado global (páginas, linhas, caracteres) é a única bagagem contínua. 

\begin{successbox}[Por que escolher o fluxo contínuo?]
Em programação competitiva, problemas com alto volume de entrada sempre priorizam fluxos iterativos, pois minimizam gargalos de largura de banda na memória cache e operam com complexidade espacial próxima de $\mathcal{O}(1)$. Soluções maduras para este problema implementam um laço dependente de \textit{tokens} padrão que consome uma palavra de cada vez, afere apenas seu atributo primordial (o tamanho) e passa para o próximo token.
\end{successbox}

\newpage

\section{Fluxograma da Solução}

Abaixo, descrevemos visualmente a simulação adotada pelo Algoritmo Guloso usando a nossa máquina de estados.

\begin{center}
\begin{tikzpicture}[node distance=1cm, auto]
  \node[wokstart] (start) {Início: $N$, $L_{max}$, $C_{max}$ \\ $P_{atual}=1$, $L_{atual}=1$, $C_{atual}=0$};
  
  \node[wokdecision, below=of start, yshift=-0.4cm] (hasword) {Há palavras a ler?};
  
  \node[wokend, right=of hasword, xshift=0.8cm] (end) {Imprime $P_{atual}$ \\ Fim};
  
  \node[wokstep, below=of hasword, yshift=-0.4cm] (read) {Lê palavra de tamanho $W$};
  
  \node[wokdecision, below=of read, yshift=-0.4cm] (firstword) {$C_{atual} == 0$?};
  
  \node[wokstep, left=of firstword, xshift=-0.3cm] (calc1) {Ocupação: $S = W$};
  \node[wokstep, right=of firstword, xshift=0.3cm] (calc2) {Ocupação: $S = W + 1$};
  
  \node[wokdecision, below=of firstword, yshift=-0.8cm] (fits) {$C_{atual} + S \leq C_{max}$?};
  
  \node[wokstep, left=of fits, xshift=-0.3cm] (updateinline) {$C_{atual} \leftarrow C_{atual} + S$};
  
  \node[wokstep, right=of fits, xshift=0.3cm] (newline) { $L_{atual} \leftarrow L_{atual} + 1$ \\ $C_{atual} \leftarrow W$};
  
  \node[wokdecision, below=of newline, yshift=-0.4cm] (newpage) {$L_{atual} > L_{max}$?};
  
  \node[wokstep, below=of newpage, yshift=-0.4cm] (addpage) {$P_{atual} \leftarrow P_{atual} + 1$ \\ $L_{atual} \leftarrow 1$};
  
  % Paths
  \draw[wokarrow] (start) -- (hasword);
  
  \draw[woknoarrow] (hasword) -- node[above] {Não} (end);
  \draw[wokyesarrow] (hasword) -- node[right] {Sim} (read);
  
  \draw[wokarrow] (read) -- (firstword);
  
  \draw[wokyesarrow] (firstword) -- node[above] {Sim} (calc1);
  \draw[woknoarrow] (firstword) -- node[above] {Não} (calc2);
  
  \draw[wokarrow] (calc1) |- (fits);
  \draw[wokarrow] (calc2) |- (fits);
  
  \draw[wokyesarrow] (fits) -- node[above] {Sim} (updateinline);
  \draw[woknoarrow] (fits) -- node[above] {Não} (newline);
  
  \draw[wokarrow] (newline) -- (newpage);
  
  \draw[wokyesarrow] (newpage) -- node[right] {Sim} (addpage);
  
  % Coordinates to route paths back to 'hasword'
  \coordinate (loop_up_left) at ([xshift=-0.6cm]updateinline.west);
  \coordinate (loop_up_right) at ([xshift=0.6cm]addpage.east);
  \coordinate (loop_up_right2) at ([xshift=0.6cm]newpage.east);
  
  \draw[wokarrow] (updateinline.west) -- (loop_up_left) |- (hasword.west);
  
  \draw[woknoarrow] (newpage.east) -- node[above] {Não} (loop_up_right2) |- (hasword.east);
  \draw[wokarrow] (addpage.east) -- (loop_up_right) |- (hasword.east);

\end{tikzpicture}
\end{center}

\newpage

\section{Análise de Complexidade}

A eficiência da solução varia sutilmente dependendo de como as strings são manipuladas. A estratégia gulosa exige unicamente uma varredura sequencial (da esquerda para a direita) pelo corpus da entrada. O processamento individual de uma palavra é constante dadas operações lógicas triviais. O que muda de fato é a alocação de memória e o tempo de construção de buffers temporários para \textit{parsing} por debaixo dos panos.

Assumiremos que existam $N$ palavras no texto total e que o comprimento total de caracteres em todo o texto (incluindo separadores) seja denotado por $\Sigma S$.

\begin{table}[h!]
\centering
\begin{tabularx}{\textwidth}{l X c c}
\toprule
\textbf{Etapa} & \textbf{Descrição} & \textbf{Tempo} & \textbf{Espaço} \\
\midrule
Abordagem Streaming & O analisador léxico extrai \textit{tokens} on-the-fly a partir do \textit{standard input stream}. Somente o tamanho de cada palavra lida é mantido. & $\mathcal{O}(N)$ ou $\mathcal{O}(\Sigma S)$ & $\mathcal{O}(\max w_i)$ \\
\midrule
Abordagem Buffering & O programa aloca contiguamente toda a string e aplica separação léxica resultando num array de referências. & $\mathcal{O}(\Sigma S)$ & $\mathcal{O}(\Sigma S)$ \\
\midrule
Lógica e Transição & Custo algorítmico do gerenciamento da máquina de estados por palavra lida (aritmética e branches). & $\mathcal{O}(1)$ por palavra & $\mathcal{O}(1)$ \\
\bottomrule
\end{tabularx}
\caption{Comparação das metodologias de leitura.}
\end{table}

A simulação gulosa aliada ao processamento em fluxo é absolutamente a abordagem de ponta e performa em tempos consistentemente perto de 0.00s e consumo de memória da ordem de kilobytes nos ambientes de juiz online, operando perfeitamente bem aos ditames do limite espacial especificado no enunciado de $195$ MB.

\vspace{1.5em}

\begin{notebox}[Tutorial Recomendado e Leituras Suplementares]
Para quem deseja aprofundar os estudos em \textbf{Algoritmos Gulosos} na aplicação de problemas em grafos e processamento de dados, sugerimos a obra literária fundamental: 
\begin{itemize}
    \item \textit{Introduction to Algorithms (Cormen, T. H. et al.)} — O Capítulo 16 descreve meticulosamente e formaliza as provas de matróides e teoremas subjacentes aos algoritmos Greedy.
\end{itemize}
Além do mais, compreender a distinção entre a versão simples que solucionamos neste editorial contra o famigerado \textbf{Text Justification Problem} ajudará incrivelmente na lapidação do seu raciocínio computacional:
\begin{itemize}
    \item \href{https://en.wikipedia.org/wiki/Line_wrap_and_word_wrap}{Wikipedia: Line wrap and word wrap}
\end{itemize}
\end{notebox}

\vfill
\wokfooter{0005}{Mundo do Código}

\end{document}
